% Image credits for the photographs and AI illustrations of Book 3
% (University Chemistry, Year 2). Machine-readable license records live in
% images/book3/CREDITS.md; keep the two files in sync. The Book 3 writer
% owns this file: add one \item per photograph (in an itemize, after the
% first paragraph) as photographs land.
% Arabic edition: personal names are transliterated as in the captions, with
% the original spelling in parentheses (the legally required credit); user
% names, institutions, works and licences stay verbatim (as in the Arabic
% Book 2 credits page, image-credits-book2.ar.tex).
\chapter*{حقوق الصور}

رسوم هذا الكتاب أصلية. أما الصور الفوتوغرافية، حيث استُعملت، فمستعملة
بمقتضى رخصها الحرة، مع الشكر لأصحابها؛ وروابط المصادر الكاملة وروابط
الرخص لكل صورة مدرجة في
\texttt{images/book3/CREDITS.md} في مستودع الكتاب.
والرسوم التحذيرية للأخطار هي رسوم النظام المنسَّق عالميًا لتصنيف
المواد الكيميائية ووسمها، كما رُسمت لحزمة
\texttt{ghsystem} في \TeX\ Live.

\bigskip
وإلى جانب الصور الفوتوغرافية والرسوم الأصلية، بعض الأشكال
رسوم مولَّدة بالذكاء الاصطناعي، أُنتجت بأداة توليد الصور من OpenAI انطلاقًا من
أوامر كتبها محررو الكتاب، ورُوجعت من حيث الدقة الكيميائية قبل
إدراجها. والأمر الكامل لكل صورة مولَّدة
مدوَّن في \texttt{images/book3/ai/PROMPTS.md} في المستودع.

\bigskip
\noindent الصور الفوتوغرافية، حسب الفصل:
\begin{itemize}\small
  \item الفصل~1: جرمان هنري هس (Germain Henri Hess)، لوحة بريشة إ.~إ. رايمرز (I.~I. Reimers)
        (1837--1838)، CC BY-SA 4.0؛ ومارسلان برتلو (Marcellin Berthelot)، مؤلف مجهول، ملكية عامة
        (كلتاهما من Wikimedia Commons).
  \item الفصل~2: جوزايا ويلارد غيبس (Josiah Willard Gibbs)، مصوّر مجهول، ملكية عامة
        (Wikimedia Commons).
  \item الفصل~3: فرانسوا ماري راؤول (François-Marie Raoult)، مصوّر مجهول، ملكية عامة
        (Wikimedia Commons).
  \item الفصل~4: ياكوبوس هنريكوس فانت هوف (Jacobus Henricus van 't Hoff)، بعدسة نيكولا پيرشايد
        (Nicola Perscheid) (1904)، ملكية عامة؛ وهنري لو شاتلييه (Henry Le Chatelier)، مؤلف مجهول،
        CC BY-SA 4.0؛ وفريتس هابر (Fritz Haber) وكارل بوش (Carl Bosch)، Nobel Foundation، ملكية عامة
        (كلها من Wikimedia Commons).
  \item الفصل~6: الفرن العالي في دويسبورغ-نورد، التقطها ديتمار رابيش (Dietmar Rabich)، CC BY-SA
        4.0؛ ولحام سكة بالثرميت، التقطها PetrS.، CC BY-SA 3.0 (كلتاهما من
        Wikimedia Commons).
  \item الفصل~7: عمود تقطير في مصفاة، التقطها CEphoto، أوفه أراناس (Uwe Aranas)،
        CC BY-SA 3.0 (Wikimedia Commons).
  \item الفصل~9: مايكل فاراداي (Michael Faraday)، بريشة توماس فيليبس (Thomas Phillips) (1842)، ملكية عامة
        (Wikimedia Commons).
  \item الفصل~10: ياروسلاف هيروفسكي (Jaroslav Heyrovsk\'y)، أرشيف معهد J.~Heyrovsk\'y
        للكيمياء الفيزيائية، CC BY-SA 3.0 (Wikimedia Commons).
  \item الفصل~11: عمود فولتا، التقطه GuidoB، CC BY-SA 3.0؛ وتشارلز مارتن هول (Charles Martin Hall)
        وبول هيرو (Paul H\'eroult)، مؤلفان مجهولان، ملكية عامة (كلها من Wikimedia
        Commons).
  \item الفصل~12: تمثال من النحاس، التقطه Elcobbola، ملكية عامة؛ ومصعد مضحًّى،
        التقطه Zwergelstern، CC BY-SA 3.0 (كلاهما من Wikimedia Commons).
  \item الفصل~13: إرفين شرودنغر (Erwin Schr\"odinger)، Nobel Foundation (1933)، ملكية عامة
        (Wikimedia Commons).
  \item الفصل~14: الأكسجين السائل في مغناطيس، التقطها بيتر كاوبر (Pieter Kuiper)، ملكية عامة؛
        وروبرت موليكن (Robert Mulliken) وفريدريش هوند (Friedrich Hund)، GFHund، CC BY 3.0 (كلتاهما من Wikimedia
        Commons).
  \item الفصل~16: أوغست كيكوليه (August Kekul\'e) (1873)، مؤلف مجهول، ملكية عامة
        (Wikimedia Commons).
  \item الفصل~18: محاليل فلزات انتقالية، التقطها Benjah-bmm27، ملكية عامة؛ وألفريد
        فيرنر (Alfred Werner)، Les Prix Nobel (1914)، ملكية عامة (كلتاهما من Wikimedia Commons).
  \item الفصل~19: بلورة ياقوت، ملكية عامة؛ وزمردة، التقطها ديدييه ديكوان (Didier Descouens)،
        CC BY-SA 3.0؛ وبلورات كبريتات النحاس، التقطها Stephanb، CC BY-SA 3.0 (كلها من
        Wikimedia Commons).
  \item الفصل~20: أكيرا سوزوكي (Akira Suzuki) وإييتشي نيغيشي (Ei-ichi Negishi) وريتشارد هيك (Richard Heck)،
        التقطها BloodIce، CC BY-SA 4.0 (Wikimedia Commons).
  \item الفصل~22: شارل فريدل (Charles Friedel) (تسعينيات القرن التاسع عشر) وجيمس ميسون كرافتس
        (James Mason Crafts) (\textit{Popular Science Monthly}، 1900)، ملكية عامة (Wikimedia Commons).
  \item الفصل~24: ميشيل أوجين شفرول (Michel Eug\`ene Chevreul)، بريشة نيكولا أوستاش موران
        (Nicolas Eustache Maurin)، ملكية عامة (Wikimedia Commons).
  \item الفصل~25: ألكسندر بورودين (Alexander Borodin)، مؤلف مجهول، ملكية عامة؛ ولودفيغ كلايزن
        (Ludwig Claisen)، التقطها Veronica.villagrasa، CC BY-SA 4.0 (كلتاهما من Wikimedia Commons).
  \item الفصل~26: روبرت روبنسون (Robert Robinson)، مؤلف مجهول، ملكية عامة (Wikimedia Commons).
  \item الفصل~28: إلياس جيمس كوري (Elias James Corey)، التقطها Trvthchem، CC0 (Wikimedia Commons).
  \item الفصل~29: والاس كاروذرز (Wallace Carothers) في المختبر، مصوّر مجهول، ملكية عامة
        (Wikimedia Commons).
  \item الفصل~30: إميل فيشر (Emil Fischer)، Atelier Victoria، برلين، نحو 1895، ملكية عامة
        (Wikimedia Commons).
  \item الفصل~31: ميخائيل تسفيت (Mikhail Tsvet)، مؤلف مجهول، ملكية عامة (Wikimedia Commons).
  \item الفصل~32: فرانسيس ويليام أستون (Francis William Aston) (1922)، مؤلف مجهول، ملكية عامة؛
        وغوستاف كيرشوف (Gustav Kirchhoff) وروبرت بنزن (Robert Bunsen) وهنري روسكو (Henry Roscoe) (1862)،
        ملكية عامة (كلتاهما من Wikimedia Commons).
  \item الفصل~34: وليام سيلي غوسيت (William Sealy Gosset) (1908)، ملكية عامة (Wikimedia Commons).
%% next chapter here
\end{itemize}
\clearpage
